% Regression: the labelled cup word draws.  `west={cup=$m$}` is the documented
% spelling for a fixed-point matrix standing on a bend; before issue 5482 the
% side key admitted only the bare alphabet, so the labelled form was rejected
% outright and every consumer hand-placed a bead of its own.  The bead the
% word mints stands at the apex of the arc, not on the chord between the two
% row ends the arc leaves.
\documentclass{standalone}
\usepackage{tenkz}
\begin{document}
\tenkzkernel
% opposite labelled sides: two distinct cup wires, each carrying its matrix
\begin{tenkz}[rows={wire,wire}, cols=3, bonds=none,
    west={cup=$\rho^R$}, east={cup=$\rho^L$}]
  \tn{} & \tn{} & \tn{}\\
  \tn[conjugate]{} & \tn[conjugate]{} & \tn[conjugate]{}
  \tnwire{(1,1)}{(1,2)} \tnwire{(1,2)}{(1,3)}
  \tnwire{(2,1)}{(2,2)} \tnwire{(2,2)}{(2,3)}
\end{tenkz}
\quad
% one labelled side beside a bare one: only the labelled side grows a bead
\begin{tenkz}[rows={wire,wire}, cols=2, bonds=none, west={cup=$X$}, east=cup]
  \tn{} & \tn{}\\
  \tn{} & \tn{}
  \tnwire{(1,1)}{(1,2)} \tnwire{(2,1)}{(2,2)}
\end{tenkz}
\quad
% the horizontal sides bend the same way about their column pairs
\begin{tenkz}[rows={wire,wire}, cols=2, bonds=none,
    west=none, east=none, north={cup=$Y$}]
  \tn{} & \tn{}\\
  \tn{} & \tn{}
  \tnwire{(1,1)}{(2,1)} \tnwire{(1,2)}{(2,2)}
\end{tenkz}
\end{document}
